\documentclass[11pt]{article}
\usepackage[T1]{fontenc}
\usepackage[utf8]{inputenc}
\usepackage{lmodern}
\usepackage{geometry}
\usepackage{enumitem}
\usepackage{hyperref}
\usepackage{fancyhdr}

\geometry{letterpaper, margin=0.9in}
\setlength{\parindent}{0pt}
\setlist[itemize]{leftmargin=1.5em, itemsep=1pt, topsep=3pt, parsep=0pt}

\hypersetup{colorlinks=true, linkcolor=blue, urlcolor=blue}

\pagestyle{fancy}
\fancyhf{}
\renewcommand{\headrulewidth}{0pt}
\fancyfoot[C]{\small\thepage}

\newcommand{\sectiontitle}[1]{\vspace{10pt}{\textbf{\large #1}}\par\vspace{-7pt}\hrulefill\par\vspace{4pt}}
\newcommand{\entry}[3]{\textbf{#1}\hfill\textit{#2 -- #3}\par}
\newcommand{\org}[1]{\textit{#1}\par}
\newcommand{\entrygap}{\vspace{6pt}}

\begin{document}

\begin{center}
    \textbf{\huge Daniel Gause} \\[6pt]
    National Radio Astronomy Observatory, Charlottesville, VA \\[3pt]
    \href{mailto:dgause@nrao.edu}{dgause@nrao.edu} $\vert$
    \href{mailto:danpgause@gmail.com}{danpgause@gmail.com} $\vert$
    \href{https://dangause.github.io}{dangause.github.io} \\[3pt]
    \href{https://github.com/dangause}{GitHub} $\vert$
    \href{https://www.linkedin.com/in/daniel-gause-b48633174/}{LinkedIn} $\vert$
    ORCID \href{https://orcid.org/0009-0004-8062-3810}{0009-0004-8062-3810}
\end{center}

\sectiontitle{Education}

\entry{B.A. Physics and Mathematics}{September 2016}{May 2020}
\org{Middlebury College}
\begin{itemize}
    \item GPA 3.5, \textit{cum laude}
    \item Physics thesis: Structural Analysis of Potential Dual Quasar Pairs
    \item Mathematics thesis: An Introduction to Large Deviation Theory
\end{itemize}

\sectiontitle{Research Experience}

\entry{Postbaccalaureate Researcher}{September 2026}{Present}
\org{National Radio Astronomy Observatory --- Advisor: Adele Plunkett}
\begin{itemize}
    \item Conducting a multiwavelength study of protostellar outflows using archival data from ALMA, the VLA, and the Hubble Space Telescope
    \item Continuing development of MANNA, with a deployment on the NOIRLab Astro Data Lab in internal testing
\end{itemize}

\entrygap
\entry{Research Intern}{May 2026}{August 2026}
\org{NSF-Simons AI Institute for Cosmic Origins (CosmicAI) --- National Radio Astronomy Observatory}
\begin{itemize}
    \item Built \href{https://github.com/dangause/manna}{MANNA}, an open-source MCP server that exposes the NOIRLab and NRAO/ALMA data archives to natural-language queries through IVOA standards (TAP, SIA, Cone Search)
    \item Served a self-hosted open-weights LLM (vLLM) on local GPUs and integrated it into the observatory's notebook platform (jupyter-ai)
\end{itemize}

\entrygap
\entry{Lead Developer, STIPS}{December 2024}{Present}
\org{University of California Observatories --- Lick Observatory}
\begin{itemize}
    \item First author and lead developer of \href{https://github.com/dangause/stips}{STIPS}, the Small Telescope Image Processing Suite, an open-source pipeline that wraps the Rubin/LSST Science Pipelines for 1-m telescopes: instrument signature removal, difference imaging, coadd-template generation, and forced photometry
    \item Validated against published results: reproduced a Nickel supernova lightcurve and calibrated Landolt standards to better than 0.05 mag with sub-pixel astrometry
    \item Extended the framework to a second instrument, the CTIO 1.0-m; runs locally, in Docker, and on Slurm/HTCondor clusters
\end{itemize}

\entrygap
\entry{Research Assistant}{May 2019}{March 2020}
\org{Department of Physics, Middlebury College}
\begin{itemize}
    \item Built Python pipelines to retrieve and process archival Hubble Space Telescope data via MAST, finding 21 candidate dual-quasar systems
    \item For senior thesis, modeled host galaxy morphologies of the three most promising candidates with GALFIT and TinyTim PSFs
\end{itemize}

\entrygap
\entry{Cataclysmic Variables Research Project}{August 2024}{2025}
\org{City College of San Francisco}
\begin{itemize}
    \item Observational project led by Professor Claia Bryja monitoring cataclysmic variable stars; collected and reduced data from the Nickel Telescope at Lick Observatory
\end{itemize}

\sectiontitle{Presentations}

\begin{itemize}
    \item Gause, D., Telkamp, Z., Plunkett, A., Mason, B., Nikutta, R., Juneau, S., Zacharia Anil, A., Loomis, R., Offner, S., Durrett, G., Torrey, P., \& Murphy, E. ``MANNA: MCP Architecture for NRAO and NOIRLab Archives.'' Poster, CosmicAI Cosmic Horizons Conference, Charlottesville, VA, July 2026.
    \item Gause, D. \& Westfall, K. ``The Small Telescope Image Processing Suite (STIPS): Bringing the LSST Stack to 1-m Telescopes.'' Poster, 248th Meeting of the American Astronomical Society, Pasadena, CA, June 2026.
    \item Gause, D. ``On the Hunt for Dual Quasars.'' Keck Northeast Astronomy Consortium Symposium, Vassar College, October 2019.
    \item Gause, D. ``Searching for DQSOs in Archival Hubble Data.'' Poster, Middlebury Summer Research Poster Session, July 2019.
\end{itemize}

\sectiontitle{Professional Experience}

\entry{Full Stack AI/ML Engineer}{May 2021}{September 2026}
\org{Accenture Federal Services}
\begin{itemize}
    \item Lead developer for a small team, owning architecture and code quality on production systems
    \item Built data and machine-learning pipelines over large geospatial datasets, including models trained on NASA satellite imagery
\end{itemize}

\sectiontitle{Volunteer and Outreach}

\entry{Scientific Computing Department Volunteer}{March 2024}{September 2026}
\org{California Academy of Sciences}
\begin{itemize}
    \item Built and deployed a Python package using machine-learning clustering for collections data
    \item Sped up a processing application by 40\% through parallelization
\end{itemize}

\entrygap
\entry{Service Learner Volunteer}{August 2024}{January 2025}
\org{City College of San Francisco}
\begin{itemize}
    \item Helped plan and host monthly public sky viewing nights and planetarium shows
\end{itemize}

\sectiontitle{Skills and Certifications}

\begin{itemize}
    \item \textbf{Languages:} Python, C++, C, SQL/ADQL, Bash
    \item \textbf{Astronomy:} LSST Science Pipelines, Gen3 Butler, Astropy, pyvo, FITS, IVOA protocols
    \item \textbf{Computing:} Slurm, HTCondor, Docker, Apptainer, HPC, CI/CD, pytest, Git
    \item \textbf{Certifications:} CompTIA Security+ (2025)
\end{itemize}

\end{document}
